% Контрпример из тетради: характеристики u_x = 0 -- прямые y = const,
% кривая y = x^3 касается одной из них в начале координат.
\begin{figure}[htbp]
\centering
\begin{tikzpicture}[>={Latex[length=4pt]}]
\begin{axis}[
  width=8.4cm, height=5.4cm,
  axis lines=middle,
  xmin=-1.75, xmax=2.3, ymin=-1.5, ymax=1.5,
  xlabel={$x$}, ylabel={$y$},
  xlabel style={anchor=north west}, ylabel style={anchor=south east},
  xtick=\empty, ytick=\empty, clip=false,
]
  % характеристики: y = const
  \foreach \c in {-1.15, -0.75, -0.35, 0, 0.35, 0.75, 1.15} {
    \addplot[blue!55!black, thick, domain=-1.6:1.9, samples=2] {\c};
  }
  % кривая y = x^3
  \addplot[very thick, red!70!black, domain=-1.12:1.12, samples=80] {x^3};

  \fill (axis cs:0,0) circle (1.8pt);
  \draw[->, gray] (axis cs:-1.1,0.95) -- (axis cs:-0.15,0.1);
  \node[anchor=south east, font=\footnotesize] at (axis cs:-1.05,0.95)
    {точка касания};
  \node[red!70!black, anchor=west, font=\footnotesize] at (axis cs:1.15,1.25)
    {$y = x^3$};
  \node[blue!55!black, anchor=west, font=\footnotesize] at (axis cs:1.95,0.35)
    {$y = \mathrm{const}$};
\end{axis}
\end{tikzpicture}
\caption{Нарушение трансверсальности. Характеристики уравнения $u_x = 0$ —
горизонтальные прямые, а кривая $y = x^3$ касается одной из них в начале
координат: при $x=0$ её касательная $(1, 3x^2)$ совпадает с $\bar a = (1,0)$.
Решение $u = y^{2/3}$ существует, но негладко на оси $Ox$.}
\label{fig:tangency}
\end{figure}
